\documentclass[a4paper,11pt]{article}
\usepackage[utf8]{inputenc}
\usepackage[T1]{fontenc}
\usepackage[french]{babel}
\usepackage[left=2cm,right=2cm,top=2cm,bottom=2cm]{geometry}
\usepackage{xcolor}
\usepackage{hyperref}
\usepackage{fontawesome5}
\usepackage{parskip}

% --- Couleurs EY ---
\definecolor{primary}{RGB}{46, 46, 56}
\hypersetup{colorlinks=true, urlcolor=primary}

\begin{document}

% --- EXPÉDITEUR ---
\noindent
\begin{minipage}[t]{0.5\textwidth}
    \textbf{Loïc Aron MBASSI EWOLO}\\
    Élève-Ingénieur Bac+5 -- Double diplôme ENSEA\\[3pt]
    \faMapMarker\ Cergy, France\\
    \faPhone\ (+33) 7 59 52 33 98\\
    \faEnvelope\ \href{mailto:loicaron.mbassiewolo@ensea.fr}{loicaron.mbassiewolo@ensea.fr}
\end{minipage}
\hfill
% --- DATE ---
\begin{minipage}[t]{0.4\textwidth}
    \raggedleft
    Cergy, le 10 mars 2026
\end{minipage}

\vspace{1.2cm}

% --- DESTINATAIRE ---
\hfill
\begin{minipage}[t]{0.5\textwidth}
    \raggedleft
    \textbf{EY}\\
    Service Recrutement\\
    Paris La Défense, France
\end{minipage}

\vspace{1cm}

\noindent\textbf{Objet :} Candidature en alternance -- Audit des Systèmes d'Information

\vspace{0.8cm}

Madame, Monsieur,

Évaluer automatiquement si une application mobile respecte ses engagements de confidentialité en croisant analyse NLP de ses politiques, analyse statique de ses SDK et permissions réellement demandées : c'est le pipeline que je développe en stage à INRIA Saclay (équipe TRIBE, sécurité et privacy). Ce travail d'analyse de conformité, qui produit un score de confiance interprétable, m'a donné une compréhension concrète de ce que signifie auditer un système d'information sous l'angle des risques, de la transparence et de la réglementation (RGPD).

Avant INRIA, deux expériences m'ont confronté aux enjeux de sécurité et de conformité des SI en environnement professionnel. Chez CSC (fintech), j'ai conçu le backend d'une plateforme de transactions bancaires : gestion de la concurrence, conformité aux normes financières, architecture sécurisée (Docker, Redis). Chez SOSDOCTA (télémédecine, 2 ans), j'ai développé et maintenu une plateforme manipulant des données médicales sensibles avec authentification OAuth2/JWT et contraintes RGPD. Ces missions m'ont appris à documenter les processus techniques, à évaluer les dispositifs de sécurité en place et à formuler des recommandations d'amélioration.

Ma maîtrise des technologies émergentes (IA générative, systèmes multi-agents, IoT) complète ce profil orienté audit SI. Avoir conçu un système de 4 agents IA (RAG, LangChain, Gemini) et travaillé sur de l'IA embarquée me permet de comprendre les risques spécifiques que ces technologies introduisent dans les processus métiers, un enjeu central de vos missions d'accompagnement. Mon passage au MILA (Institut Québécois d'IA) en environnement anglophone et mes rédactions scientifiques en anglais à INRIA garantissent mon aisance dans les contextes internationaux qu'EY propose.

Je suis disponible dès septembre 2026 pour une alternance de 12 mois. Cartographier les risques SI de clients dans des secteurs variés (luxe, aéronautique, énergie), tout en développant une expertise en audit et certification (ISO 27001, SOC, ISAE 3402), est le type de mission dans lequel je souhaite m'investir.

Je reste à votre disposition pour un entretien afin de discuter de ma candidature.

Je vous prie d'agréer, Madame, Monsieur, l'expression de mes salutations distinguées.

\vspace{0.8cm}

\textbf{Loïc Aron MBASSI EWOLO}\\
\textit{Élève-Ingénieur ENSEA / Polytechnique Yaoundé}

\end{document}
